\section{Pilot Games and Learning: Verified Play and Local Education}

Pilot Fabric 0.46.0 completes the locally enforceable portion of master-list
Section~8.  The release treats cloud gaming and children's learning as two governed
experiences rather than as links to third-party pages.  Both experiences remain useful
without a paid language-model API, and both distinguish observed facts from assumptions.

\subsection{Governed cloud-gaming state machine}

The \texttt{GovernedGamingExperience} service persists a bounded local record for each
gaming attempt.  Its principal states are
\begin{center}
\texttt{provider\_surface\_prepared} $\rightarrow$
\texttt{provider\_surface\_observed} $\rightarrow$
\texttt{search\_results\_observed} $\rightarrow$
\texttt{title\_observed} $\rightarrow$
\texttt{stream\_ready\_observed} $\rightarrow$
\texttt{playing\_verified}.
\end{center}
Opening GeForce NOW therefore does not prove that its page appeared; a search does not
prove that a game launched; and a game title does not prove that a stream is ready or
playing.  The final playing state requires signed provider or device evidence after a
verified title and stream-readiness observation.  This prevents the television UI or a
language model from manufacturing a successful-play claim from URL delivery alone.

Provider sign-in remains provider-controlled.  Pilot opens only the allowlisted official
origin, prefers QR or device-code handoff, never stores a password in gaming state, and
never exposes credentials to a device node.  CAPTCHA, terms, multi-factor authentication
and the final sign-in decision remain owner actions.  The existing private keyboard is
ephemeral; a real password must be entered only through the trusted HTTPS controller or
the provider's own device flow.

\subsection{Controller qualification}

The private phone surface uses the browser Gamepad API, when available, to inspect
advertised axes and buttons.  Pilot maps these observations onto bounded profiles:
standard gamepad, LG mobile gamepad, keyboard and pointer, navigation-only TV remote, or
unknown.  Advertisement alone produces
\texttt{advertised\_not\_field\_verified}.  A compatible input mode is marked verified
only after an actual button or axis event is observed.  Navigation-only devices are not
reported as gameplay-ready.

Controller identifiers are hashed in the audit record.  Raw button histories are not
used as behavioural profiles.  A device can be re-qualified after browser, television,
controller or firmware changes.

\subsection{Private game shortcuts and continuation}

A saved-game shortcut can be created only from a provider title that has been independently
observed.  The shortcut is persona-bound and stores the provider, verified title and an
optional provider content identifier.  It contains no provider credential or entitlement.
The command ``Pilot, continue my game'' prepares the last private shortcut and reopens the
official provider surface; it deliberately reports \texttt{prepared\_not\_playing} until
fresh readiness and playback evidence arrive.  Different household personas do not see or
continue one another's private shortcuts.

\subsection{Expanded local curriculum}

The learning engine now uses schema \texttt{aion.education.v2}.  It supports the age bands
6--7, 8--9 and 10--12; foundation, developing and challenge difficulty; and Spanish,
mathematics and science subjects.  Spanish includes a larger word set, adaptive weak-item
selection and deeper conversation rounds.  Mathematics and science use deterministic,
offline activities that do not require web access or paid inference.

Each round contains an original code-generated SVG illustration with alternative text.
These local vectors replace emoji-only teaching and introduce no third-party copyright or
tracking dependency.  Correctness is expressed with text, symbols and colour together.
The television interface supports directional focus, visible prompts, spoken prompts and
reduced-motion rendering.

\subsection{Local pronunciation and child-safe retention}

Pronunciation scoring consumes an on-device speech transcript, normalizes case and
diacritics, and combines sequence similarity with expected-word coverage.  The child sees
an encouraging score band and retry guidance.  Pilot retains only the card identifier,
numeric score, band and time.  Raw audio and recognized words are discarded and are not
written to the learner file or audit ledger.  No cloud model is required.

\subsection{Guardian-private progress and corrections}

Progress reports are available only on the persona-bound private phone.  They show
attempts, accuracy, pronunciation average, curriculum settings and low-mastery activities.
A guardian can attach a bounded practice correction; the record contains a one-way persona
hash rather than the private identifier.  Shared-TV speech cannot inspect or alter these
reports.  Learner profiles are local aliases, not child accounts, advertising identities
or externally addressable profiles.

\subsection{Safety and assurance boundary}

The local technical review verifies keyboard and directional focus, visible and spoken
prompts, colour-independent feedback, disabled open chat, disabled external links,
disabled purchases, no adverts and no child raw-voice retention.  This does not substitute
for independent pedagogy, accessibility and child-safety assessment.  Those consumer
claims, together with provider telemetry and hardware-matrix field qualification, remain
in the Section~8 blocker register.

\subsection{Verification}

Automated checks cover state separation, signed-playing requirements, controller
qualification, persona isolation, shortcut preconditions, subject and difficulty routing,
original vector output, transcript non-retention and guardian authorization.  The full
Fabric suite must pass before the release archive is built.  Tests never claim a real
GeForce stream, a real controller model or an externally reviewed curriculum without the
corresponding field evidence.

